\chapter{1968:Lars Onsager}

\label{chap:chemistry1968}

The Nobel Prize in Chemistry for the year 1968 was awarded to Lars Onsager.

\textbf{Motivation:}

for the discovery of the reciprocal relations bearing his name, which are fundamental for the thermodynamics of irreversible processes

\section{Lars Onsager}

\subsection*{Early Life and Education}

Lars Onsager was born on 1903-11-27 in Oslo, Norway.

\subsection*{Career and Major Research}

Onsager studied chemical engineering at the Norwegian Institute of Technology in Trondheim. After positions in Europe and the United States, he joined Yale University, where he was professor of theoretical chemistry. He derived the reciprocal relations of irreversible thermodynamics and later solved the two-dimensional Ising model exactly.

\subsection*{Nobel Prize-winning Work}

The work for which Lars Onsager was awarded the Nobel Prize in Chemistry in 1968 was described by the Nobel Committee as: "for the discovery of the reciprocal relations bearing his name, which are fundamental for the thermodynamics of irreversible processes".

Onsager showed that the coefficients linking fluxes and forces in non-equilibrium systems obey symmetry relations, providing a general foundation for the thermodynamics of irreversible processes.

\subsection*{Significance and Impact}

The Onsager reciprocal relations underlie the description of transport phenomena such as diffusion, heat conduction, and thermoelectric effects and are fundamental to non-equilibrium thermodynamics.

\subsection*{Later Career and Honors}

Onsager remained at Yale until 1972, then held a research professorship at the University of Miami in Coral Gables, Florida. He died on 1976-10-05 in Coral Gables, Florida, USA.

\subsection*{Legacy}

Onsager's reciprocal relations and his exact solution of the Ising model are landmarks of twentieth-century statistical mechanics.

\subsection*{Selected Publications}

"Reciprocal Relations in Irreversible Processes" (Physical Review, 1931)

"Crystal Statistics. I. A Two-Dimensional Model with an Order-Disorder Transition" (Physical Review, 1944)

\clearpage

\section*{Chapter Summary}

The 1968 prize honored Lars Onsager for the discovery of the reciprocal relations of irreversible thermodynamics. His work provides the foundation for understanding transport processes in chemistry and physics.

